% section: 漸化式の分類

\section{漸化式の分類}
  具体的な解法について触れる前に,解法の分類を行おう.
  まず漸化式は解ける漸化式と解けない漸化式に分類される.
  そして解ける漸化式は「決まった手順で解ける漸化式」と「特定の性質から解ける漸化式」に分類される.
  「決まった手順で解ける漸化式」はさらに分類され,各分類ごとに様々な解法がある.
  ここで言う「解ける」というのは,本書内での定義であり,一般的なものではない.
  ここで解けないとして扱われるような漸化式は一般項を表す手段が存在しないわけではない.
  あくまで本書で扱う手法の内では表現できないという意味である.
  「解ける」の定義について述べておく.
  定義は高校数学程度の範囲内で既知の関数に帰着するか,有限回の演算または操作により一般項を得られるような解法が存在するものを意味している.

  これらの分類を元に,まずは漸化式の基本的なパターンを基本4パターンの章で説明する.
  続く応用パターンで,いかにして基本パターンに帰着させて漸化式を解けばよいのかを議論する.

